\input{preamble}

\FirstPage{1}
\LastPage{x}

\pdfcrop
%\draft

\begin{document}

\udc{XXX.XXX}
\msc{YYYYY, ZZZZZ, DDDDD}

\rutitle{Введите название Вашей статьи здесь}
\rutitleshort{Короткое название статьи, если необходимо (eсли не помещается сверху)}

\entitle{Enter your paper title here}
\entitleshort{Enter the short paper title here if necessary}

\ruauthor{А.~А.~Иванов\affil{1,2}, В.~В.~Смирнов\affil{2}}{И\,в\,а\,н\,о\,в~А.~В., С\,м\,и\,р\,н\,о\,в~В.~В.}
\enauthor{A.~A.~Ivanov\affil{1,2}, V.~V.~Smirnov\affil{2}}{I\,v\,a\,n\,o\,v~A.~V., S\,m\,i\,r\,n\,o\,v~V.~V.}

\ruaffil{Московский государственный университет имени М.~В.~Ломоносова,\\ 
Научно-исследовательский институт механики, \\
Россия, 119192, Москва, Мичуринский проспект, 1.
\and 
Самарский государственный технический университет,\\
Россия, 443100, Самара, ул. Молодогвардейская, 244.}

\enaffil{Lomonosov Moscow State University,\\ 
Institute of Mechanics, \\
1, Michurinsky prospekt, Moscow, 119192, Russian Federation.
\and 
Samara State Technical University,\\
244, Molodogvardeyskaya st., Samara, 443100, Russian Federation.}

\ruauthorsdetails{2}{% Количество авторов
\fullauthor{Андрей Алексеевич Иванов}
\CAuthor % Автор-корреспондент
\orcid{0000-0000-0000-xxxx} % ORCID ID автора 
\regalia{кандидат физико-математических наук, доцент} 
\position{старший научный сотрудник} 
\dept{лаб. упругости и пластичности} 
\email{anton@ya-mail.ru}
\fullauthor{Владимир Викторович Смирнов}
\orcid{0000-0000-0000-xxxx} % ORCID ID автора 
\regalia{доктор физико-математических наук, профессор} 
\position{заведующий кафедрой} 
\dept{каф. упругости и пластичности} 
\email{vlad_smirnov@ya-mail.ru}
}

\enauthorsdetails{2}{
\fullauthor{Andrew A. Ivanov}
\CAuthor % Сorresponding Author
\orcid{0000-0000-0000-xxxx} % ORCID ID
\regalia{Cand. Phys. \& Math. Sci., Associate Professor} 
\position{Senior Researcher} 
\dept{Lab. of Elasticity and Plasticity}
\email{anton@ya-mail.ru}
\fullauthor{Vladimir V. Smirnov}
\orcid{0000-0000-0000-xxxx} % ORCID ID avtora 
\regalia{Dr. Phys. \& Math. Sci., Professor} 
\position{Head of Department} 
\dept{Dept. of Elasticity and Plasticity}
\email[br]{vlad_smirnov@ya-mail.ru}
}

\rukeywords{ключевое слово, ключевая фраза, новая мысль}
\enkeywords{keyword, key phrase, new thought}

\ruabstract{Аннотация на русском языке должна содержать 3–5 предложений (не более 500 знаков без пробелов). Рекомендуется структурировать аннотацию: объект исследования, цель, использованные методы и подходы, основные результаты. Аннотация должна быть информативной, отражать основное содержание статьи и результаты исследований, следовать логике описания результатов в статье и не содержать общих слов. Аннотация не должна содержать ссылок на литературу и аббревиатуры (если это возможно).}

\enabstract{An abstract is a succinct summary of a longer piece of work, usually academic in nature, which is published in isolation from the main text and should therefore stand on its own and be understandable without reference to the longer piece. It should report the latter's essential facts, and should not exaggerate or contain material that is not there. Its purpose is to act as a reference tool (for example in a library abstracting service), enabling the reader to decide whether or not to read the full text. 

The abstract should reflect only what appears in the original paper. The Abstract should:
\begin{itemize}
\item Describe the main objective(s) of the study;
\item Explain how the study was done, including any model organisms used, without methodological detail;
\item Summarize the most important results and their significance.
\end{itemize}
Abstracts should not include citations and abbreviations, if possible.}

\makerutitle



\Section[n]{Введение}
Это шаблон статьи для представления рукописей докладов на X Всероссийскую конференцию по механике деформируемого твердого тела.

\Section{Требования к статье}
Редакция просит авторов при оформлении работ придерживаться следующих правил и рекомендаций:
\begin{enumerate}
\item  Объем статьи "--- не более 3 страниц, включая список литературы, рисунки и таблицы, аннотации на русском и английском языках.
\item Статья должна содержать: название работы, список авторов, представленный в алфавитном порядке; аннотацию (объем "--- до 500 знаков без пробелов), которая дается перед основным текстом; список ключевых слов; основной текст, который рекомендуется разделять на подразделы с целью облегчения чтения работы; заключение с краткой характеристикой основных полученных результатов; название работы на английском языке с указанием всех авторов; список ключевых слов на английском языке; аннотацию на английском языке; сведения о всех авторах на русском и английском языках: должность, степень, звание, вуз, его полный почтовый адрес, email. Название работы должно адекватно отражать ее содержание и быть, по возможности, кратким. Не допускается включение формул в название работы и текст аннотации.
\item Статья должна быть снабжена индексом универсальной десятичной классификации (УДК) и классификатором MSC (\url{http://www.ams.org/msc/}).
\item Текст статьи должен быть подготовлен средствами издательской системы \LaTeXe\ с использованием стиля \href{https://www.overleaf.com/read/dknkxrfwrdsw}{\tt mdtt-samara.sty}. 
Рисунки представляются отдельно в~формате {\tt wmf}, {\tt png} с разрешением не менее 600 dpi. 
Изменение стандартных стилевых файлов недопустимо.
\item Библиографические ссылки оформляются в соответствии с ГОСТ Р 7.0.5-2008  c помощью пакета {\tt samgtu-bib.sty}. 
\item Статью в формате {\tt tex} и {\tt pdf} версию необходимо загрузить через личный кабинет участника конференции (\url{http://mdtt.su/personal/}) и продублировать по электронному адресу \href{mailto:radch@samgtu.ru}{\tt radch@samgtu.ru}. 
Для подачи также можно использовать share-ссылку из \href{https://www.overleaf.com/}{Overleaf}.

\end{enumerate}


Обратите внимание что в нашем стилевом файле используется команда секционирования \verb"\Section". 
Просим Вас  использовать именно её. 
Если раздел не надо нумеровать, то воспользуетесь вариантом \verb"\Section[N]".

\Section{Набор формул и организации ссылок}
Для набора формул Вы можете пользоваться любым руководством по \LaTeXe.
Мы рекомендуем использовать краткое руководство от К.~В.~Воронцова (\url{http://www.ccas.ru/voron/download/voron05latex.pdf}).

Ссылки на формулы оформляются только командой \verb"\eqref".
Ссылки на таблицы и рисунки оформляются командой \verb"\ref".
Для ссылок на литературу можно использовать команды \verb"\cite" и \verb"\citen".
Обратите внимание, что в нашем журнале нумеруются только те объекты (формулы, рисунки, таблицы, элементы списка литературы), на которые есть ссылки. 
Ссылки на формулы или другие объекты, которые не <<помечены>>, будут обозначаться так ({\bf ??}) или так [{\bf ?}].


\Section{Пользуйтесь окружениями для теорем, лемм, задач и прочего}
Вы можете использовать следующие определенные в пакете {\tt samgtu} окружения {\tt definition}, {\tt lemma}, {\tt example}, {\tt remark}, {\tt task}, {\tt proof}.

\begin{lemma}
Это текст леммы.
\end{lemma}

\begin{remark}[1]
Это текст замечания. Номер окружениям  {\tt definition}, {\tt lemma}, {\tt example}, {\tt remark}, {\tt task} проставляется пока вручную.
\end{remark}


\begin{newthm}{Теорема} Это теорема.
\end{newthm}

\smallskip

\begin{proof}
Это доказательство теоремы
\end{proof}

\Section{Вставляйте рисунки и таблицы} 
Рисунки и таблицы вставляются с помощью окружений {\tt figure} и {\tt table}. 
Обратите внимание, что рисунки размещаются в каталоге {\tt ./00PIC/}.
 
\begin{table}[h!]

\small \centering

\caption{Значения параметров аппроксимации и погрешности для дробного аналога модели Кельвина  \label{un:tabl1}}
\vspace{-3mm}
\renewcommand{\tabcolsep}{0.55cm}
\begin{tabular}{c||c|c|c|c|c}
\hline
\rule{0mm}{11pt}%
\footnotesize $\sigma_{0}$, MPa & \footnotesize  $\alpha$ & \footnotesize  $E_{1}$, MPa & \footnotesize  $E_{2}$, MPa & \footnotesize  $\eta$ & \footnotesize  $\Delta,\, \%$ \\
            \hline
            \rule{0mm}{12pt}%
            \phantom{1}4.655 & 0.355 & 145.525 & 1163.800 & 190.067 & 2.586 \\
             \phantom{1}6.288 & 0.326 & 133.601 & \phantom{1}898.329 & 158.793 & 4.097 \\
             \phantom{1}8.738 & 0.394 & 140.194 & 1028.000 & 128.956 & 3.168 \\
             10.372 & 0.318 & \phantom{1}97.571 & \phantom{1}829.733 & 158.840 & 3.468 \\
             12.005 & 0.400 & 127.962 & \phantom{1}857.500 & 113.856 & 3.121 \\
            \hline
        \end{tabular}

\end{table} 


В табл.~\ref{un:tabl1} приведены значения параметров аппроксимации для дробного аналога модели Кельвина. 



\begin{figure}[t!]
\centering
\includegraphics[width=0.6\textwidth]{Kelvin570g}

\caption{Экспериментальная (сплошная линия) и расчетная (штриховая линия) кривые вязкоупругого деформирования поливинилхлоридного пластиката при напряжении $\sigma_{0}=4.655$~МПа с~последующей разгрузкой
\label{Kelvin570g}
}

\end{figure}




\Section{Немного про библиографический список}
Для оформления библиографического здесь используется пакет {\tt samgtu-bib}. 
Этот пакет "--- модификация пакета {\tt amsbib}, который размещён на сайте \href{http://www.mathnet.ru/poffice/amsbibpackage.phtml?wshow=amsbibpackage&option_lang=rus}{Math-Net.Ru} и  используется в~ведущих журналах Отделения математических наук РАН. 
Поэтому для пакета {\tt samgtu-bib} верны все примеры, которые приводятся в руководстве к пакету {\tt amsbib} (\url{http://www.mathnet.ru/poffice/amsbib.pdf}).

Просим Вас приводить в библиографическом списке выверенную информацию.
Не приводите источники, которые трудно найти, например, тезисы конференций, препринты, не имеющие электронного варианта в сети интернет.
Также не рекомендуем приводить ссылки на web-ресурсы, не являющиеся научными журналами.

Старайтесь в библиографическом списке приводить оригинальную, не переводную литературу.

Если Вы ссылаетесь на статью, размещенную на портале \href{http://www.mathnet.ru/}{Math-Net.Ru}, то смело может копировать её цитирование в формате  {\tt amsbib} с~портала.


\Section[N]{Заключение (Выводы)}
Каждая статья должна заканчиваться заключением или выводами по работе. 
Приведите здесь краткую формулировку результатов своего исследования. 
После этого раздела (перед библиографическим списком) Вы можете выразить благодарности своим коллегам или предоставить другую информацию с помощью команды \verb"\AdditionalContent".

Не забудьте сгенерировать страницу с информацией о статье на английском языке с помощью команды \verb"\makeentitle".


\AdditionalContent{Благодарность}{Работа поддержана \dots .}

\begin{thebibliography}{10}

\RBibitem{RepKum16}
\by Репин~О.~А., Кумыкова~С.~К.
\paper Внутренне-краевая задача с операторами Римана--Лиувилля для уравнения смешанного типа~третьего порядка
\jour Вестн. Сам. гос. техн. ун-та. Сер. Физ.-мат. науки
\yr 2016
\vol 20
\issue 1
\pages 43--53
\crossref{10.14498/vsgtu1461}


\RBibitem{ulg:bib:Rabotnov1966}
\by Работнов~Ю.~Н.
\book Ползучесть элементов конструкций
\yr 1966
\publ Наука
\publaddr М.
\totalpages 752

\Bibitem{ulg:bib:Uchaikin2013}
\by Uchaikin~V.~V.
\paper Heredity and Nonlocality
\inbook Fractional Derivatives for Physicists and Engineers
\yr 2013
\pages 3--58
\publ Springer Berlin Heidelberg
\publaddr Berlin
\bookvol 1
\volinfo Background and Theory
\crossref{10.1007/978-3-642-33911-0_1}


\RBibitem{ulg:bib:Bronsky}
\by Бронский~А.~П.
\paper Явление последействия в твёрдом теле
\jour ПММ
\yr 1941
\vol 5
\issue 1
\pages 31--56


\end{thebibliography}

\makeentitle

\end{document}
